% Proxy–True Decoupling under Speed-Reward Distortion in CartPole PPO,
% with a Reproducible Chained-Task ETL Audit
%
% Status: internal draft (2026-10-11), NOT submitted.
% Every number in this file is copied from a landed artefact; the pointer is in
% the comment above the table or the sentence. No number is computed by hand.
% Venue target (per PAPER_DRAFT/OUTLINE.md): Safe RL / Safe Control workshop first.
\documentclass[11pt]{article}
\usepackage[margin=1in]{geometry}
\usepackage{booktabs}
\usepackage{graphicx}
\usepackage{amsmath}
\usepackage{hyperref}
\usepackage{xcolor}

\newcommand{\todo}[1]{\textcolor{red}{\textbf{[TODO: #1]}}}
\newcommand{\exploratory}{\textcolor{blue}{\textbf{[exploratory]}}}
\emergencystretch=3em
\newcommand{\tbd}[1]{\textcolor{red}{#1}}

% Figures are PDFs generated by src/_plot_style.py, copied next to main.tex.
% Cited for existence only; no value is read off any figure by eye.
\graphicspath{{../c2/results/}{../etl/results/}}

\title{Proxy--True Decoupling under Speed-Reward Distortion in CartPole PPO,\\
with a Reproducible Chained-Task ETL Audit}
\author{Anonymous submission}
\date{Draft --- 2026-10-11}

\begin{document}
\maketitle

\begin{abstract}
We measure whether a speed-shaped proxy objective decouples from the true
objective it is supposed to approximate. On CartPole we run a preregistered
grid of PPO trainings in which the per-step reward is $\mathrm{speed}^{\alpha}$
for $\alpha \in \{0.5, 1.0, 1.5, 2.0, 2.5, 3.0\}$, five seeds per level, and we
score every policy against a ground-truth objective that penalises both
falling and speeding. The true score peaks at $\alpha=0.5$
($444.02$, 95\% CI $[431.89, 456.15]$) and then collapses: $\alpha=1.0$
reaches $-2.22$ (CI $[-4.52, 0.08]$) while the proxy score at that same level
is the largest observed over the whole grid ($1.38$), so the proxy ranking and
the true ranking are inverted exactly where the proxy is maximised. Under the
preregistered decision rule this is recorded as ``decoupling confirmed'',
with the third condition's literal clause annotated as not satisfied
(its intent is met). We additionally audit a chained-task
measurement pipeline of 90 tasks and 614 steps: the judge is replay-verified
(mismatch $=0$ over 614 steps). A first batch (one local model, task rules
hidden) lands at floor (5/614 steps); we trace the floor to task design
rather than model capability, make the rules public while the truth chain
stays validator-only, and rerun: semantic correctness rises off the floor
and the error decomposition gains a support surface (cumulative errors
dominate; an error-then-recover gradient across families is compatible with
re-anchoring from the public prompt, reported as exploratory). Both the
null and the design-failure diagnosis are reported as measurements.
\end{abstract}

\section{Introduction}
\label{sec:intro}

Two failure modes recur when an agent is trained or evaluated against
something other than the quantity we actually care about. The first is
\emph{proxy failure}: the optimisation target is a proxy for the true goal,
and maximising the proxy can move the true goal in the wrong direction
\cite{goodhart1975problems,campbell1979assessing}. The second is
\emph{measurement failure}: the procedure that scores an agent is itself
buggy, so what we report is not what the agent did
\cite{strathern1997improving,sainz2024nlp}. This draft addresses both
failure modes separately and deliberately does not address the third
(unfalsifiable claims), which is out of scope for the present measurements
\cite{popper1963conjectures}.

Our contribution is a measurement discipline more than a method. For proxy
failure we fix the proxy family before touching the trainer, freeze the
decision rule, and report the full six-level curve rather than the
single best level. For measurement failure we build a judge whose verdicts
can be replayed exactly, and we report the floor effect instead of hiding it
behind an aggregate.

\paragraph{What we do not claim.} Nothing in this draft is a capability
claim. The chained-task experiment shows that the measurement pipeline is
sound and that the evaluated model performs at floor; it does not show that
any architecture is better or worse at chained reasoning.
\section{Proxy--true decoupling under speed-reward distortion}
\label{sec:decoupling}

\subsection{Design}

The true objective is a linear combination of survival and speed regulation:
\begin{equation}
\mathrm{true} = \mathrm{steps\_alive}
  - \lambda \cdot \max(0,\; |\dot{x}| - v_{\mathrm{limit}}),
\end{equation}
with $\lambda = 1.0$ and $v_{\mathrm{limit}} = 0.376$ frozen before any run
(\texttt{prereg/preregistration\_c2.md:\,L50}). The proxy objective trained
against is the per-step reward $\mathrm{speed}^{\alpha}$, with
$\mathrm{speed} = |\dot{x}|$, over the grid
$\alpha \in \{0.5, 1.0, 1.5, 2.0, 2.5, 3.0\}$, five seeds per level
(42--46), $100{,}000$ timesteps, deterministic evaluation over 50 episodes
(\texttt{c2/results/train\_config.json:\,L2--L19}).

One deviation from the preregistration is declared rather than hidden: the
agent is trained on the per-step dense reward $\mathrm{speed}^{\alpha}$, while
the reported proxy score is $\mathrm{mean}(\mathrm{speed})^{\alpha}$. The two
are not identical, and the declaration is landed at
\texttt{train\_config.json:\,L25}. We report both the true score and the proxy
score at every level, and the preregistration requires \texttt{steps\_alive}
and \texttt{mean\_speed} to be reported alongside them.

\subsection{Preregistered decision rule}

The detachment ratio is
\begin{equation}
\mathrm{detach} = 1 - \frac{\mathrm{true}}{\mathrm{true}_{\max}},
\end{equation}
and the three-way rule is: \emph{decoupling confirmed} if the true score drops
by at least 15\% after the peak \emph{and} the 95\% CI of the drop excludes
zero; \emph{partially} for a 5--15\% drop; \emph{refuted} if the true score is
monotonically non-decreasing (\texttt{prereg/preregistration\_c2.md:\,L63--L68}).
The power analysis caps the test-relevant SD at $44.6$
(\texttt{L81--L86}); the post-run paired differences are
$446.2 \pm 15.1$ and $432.7 \pm 16.0$ (\texttt{L88--L98}).

\subsection{Results}

% Source: goodhart-project/c2/results/table1_true_score.csv, transcribed verbatim.
% Raw per-seed records: c2/results/summary.csv (header + 30 rows), sha256 58b64fc8…192c31.
\begin{table}[t]
\centering
\caption{True score with 95\% CI, proxy score, and detachment ratio across the
$\alpha$ grid. Six levels, five seeds each. Values are copied verbatim from
\texttt{c2/results/table1\_true\_score.csv}.}
\label{tab:c2}
\begin{tabular}{ccccc}
\toprule
$\alpha$ & true score (mean) & 95\% CI & proxy score & detach \\
\midrule
0.5 & $444.0216$ & $[431.894, 456.149]$ & $0.6429$ & $0.0$ \\
1.0 & $-2.2205$  & $[-4.519, 0.081]$     & $1.3762$ & $1.005$ \\
1.5 & $39.4061$  & $[31.392, 46.777]$    & $1.0667$ & $0.9113$ \\
2.0 & $59.4365$  & $[47.910, 71.084]$    & $0.8549$ & $0.8661$ \\
2.5 & $21.3227$  & $[16.043, 30.021]$    & $1.0641$ & $0.952$ \\
3.0 & $11.341$   & $[10.459, 12.472]$    & $0.8108$ & $0.9745$ \\
\bottomrule
\end{tabular}
\end{table}

Measurements only:
\begin{itemize}
  \item The true score peaks at $\alpha=0.5$ ($444.02$) and falls to
        $-2.22$ at $\alpha=1.0$, then to $11.34$--$59.44$ for
        $\alpha \ge 1.5$; the post-peak drop spans $86.6\%$ to $100.5\%$.
  \item The proxy score reaches its grid-wide maximum at $\alpha=1.0$
        ($1.38$), which is exactly the level where the true score collapses
        to $-2.22$ (CI upper bound $0.081$). At that level the proxy ranking
        and the true ranking are inverted.
  \item Under the preregistered rule this batch is recorded as
        ``decoupling confirmed''. The third condition is not satisfied
        literally, only in intent; the adjudication is landed at
        \texttt{docs/HANDOVER.md:\,\S13.11} (the three-condition check and
        the final report wording are both in that subsection).
\end{itemize}

\subsection{Limitations}

Five seeds, one environment (CartPole), one reward family
($\mathrm{speed}^{\alpha}$). The measured seed SD of $15$--$16$ is below the
preregistered $44.6$ cap, so power is approximately $1.0$ for \emph{this}
criterion only; the claim does not extend to other criteria. The $\alpha^{*}$
that maximises the true score is left-censored at the grid's left edge
(every seed peaks at $\alpha=0.5$), so its variance is degenerate and no
positive statement about $\alpha^{*}$ stability is made.

\section{A replay-verifiable chained-task audit}
\label{sec:etl}

\subsection{Task families and judge}

Three families of 30 tasks each (90 total): \texttt{stateful-calc}
(stateful computation chains), \texttt{chain-sort} (chained sorting), and
\texttt{ret-path} (backtracking path search). The judge is a one-way
validator: the ground-truth schedule is read only by the validator and is
never visible to the model (\texttt{src/validators.py:\,Validator.check\_step}
scores each step against \texttt{truth\_chain[k]}). A fail-closed gate
(\texttt{src/leak\_check.py}) probes each family and fails the whole run if
any probe answers with the original ground truth more than a calibrated
noise band allows (v3.1: a leak signal must dominate the executor's own
guess-and-miss noise, both in absolute count and relative to honest
execution). Afterwards \texttt{src/audit\_replay.py} replays every step and
requires mismatch $= 0$.

Two batches are reported. Batch~1 evaluates \texttt{llama3.1:8b} at
temperature $0$ with task rules hidden from the prompt. Batch~2 evaluates
\texttt{qwen2.5:3b} at temperature $0$ after the per-step rules
(\texttt{ops}/\texttt{sort\_rule}/\texttt{ptr}/\texttt{values}) are included
in the public prompt while \texttt{truth\_chain} remains validator-only; the
schedule file hash is unchanged across batches, so the judge scores the same
ground truths. Choosing local models at temperature $0$ is a
reproducibility choice, not a claim that the configuration is best or
strongest.

\subsection{Results: a floor, its diagnosis, and the fix}

% Source: goodhart-project/etl/results/archive_20261008_hidden_rules/table_detach.csv, verbatim.
% Batch 1: llama3.1:8b, task rules hidden. step_acc = semantic correctness
% scored by Validator.check_step() against ground truth; parse_ok = JSON
% well-formedness.
\begin{table}[t]
\centering
\caption{Batch 1 (rules hidden): 90 tasks / 614 steps, \texttt{llama3.1:8b}.
\texttt{step\_acc} is semantic correctness against the hidden schedule;
\texttt{parse\_ok} is JSON well-formedness; \texttt{first\_err\_median} is the
median index of the first erroneous step.}
\label{tab:etl}
\begin{tabular}{lccc}
\toprule
family & step\_acc & parse\_ok & first\_err\_median \\
\midrule
chain-sort    & $0.0$    & $0.0$    & $0.0$ \\
ret-path      & $0.0$    & $0.0$    & $0.0$ \\
stateful-calc & $0.0235$ & $0.0235$ & $0.0$ \\
\bottomrule
\end{tabular}
\end{table}

% Source: goodhart-project/etl/results/archive_20261012_public_rules/table_detach.csv, verbatim.
% Batch 2: qwen2.5:3b, per-step rules public, truth chain validator-only.
\begin{table}[t]
\centering
\caption{Batch 2 (rules public): 90 tasks / 614 steps, \texttt{qwen2.5:3b}.
Same judge, same schedules (file hash unchanged).}
\label{tab:etl-v2}
\begin{tabular}{lccc}
\toprule
family & step\_acc & parse\_ok & first\_err\_median \\
\midrule
chain-sort    & $0.1168$ & $0.1168$ & $0.0$ \\
ret-path      & $0.1912$ & $0.1912$ & $0.0$ \\
stateful-calc & $0.1174$ & $0.1174$ & $1.0$ \\
\bottomrule
\end{tabular}
\end{table}

\begin{itemize}
  \item The judge pipeline is reproducible in both batches: replaying all
        614 steps gives mismatch $= 0$ each time (exit status $0$).
  \item Batch 1 is at floor: semantic correctness $0/197$, $0/204$, and
        $5/213$ for the three families ($5/614 = 0.8\%$ overall); no task is
        fully correct ($0/90$).
  \item The batch-1 floor is a design property, not a capability
        measurement: the public fields of each family are insufficient to
        derive the ground truth (the integer operators, the sort rule, and
        the pointer/value mapping were never in the prompt), so the tasks
        are unsolvable by construction. The gate passing at floor is a
        statement about the probe set, not about the model.
  \item Batch 2 (rules public, truth chain still validator-only) lifts the
        floor: $0.1168$/$0.1912$/$0.1174$ per family, and
        \texttt{first\_err\_median} moves to $1.0$ for \texttt{stateful-calc}
        (some tasks now survive step zero).
\end{itemize}

\subsection{Error decomposition gains a support surface after the fix}
\label{sec:etl-decomp}

\exploratory{} The single-vs-cumulative error decomposition is the most
original part of the analysis and is \emph{not} preregistered, so it is
marked exploratory at the section head. Batch 1 had a floor effect
(\texttt{first\_err\_median} $0.0$ everywhere): the model erred at step zero,
so there were no samples of an agent continuing on a corrupted state, and
the decomposition had no support surface (\texttt{docs/HANDOVER.md:\,\S13.13}).
Batch 2 removes the floor, and the decomposition becomes measurable
(\texttt{qc/q2\_q4\_analysis/results.json}): cumulative errors dominate in
all three families ($154/188 = 82\%$, $133/174 = 76\%$, $114/165 = 69\%$ for
stateful-calc, chain-sort, ret-path respectively). Error-then-recover
transitions occur in $46/614$ steps, and the recovery rate orders across
families as ret-path $27/165 = 16.4\%$, chain-sort $15/174 = 8.6\%$,
stateful-calc $4/188 = 2.1\%$. This ordering is compatible with a
re-anchoring account: after the rules became public, a model that loses the
state can recompute it from the prompt, and the ease of recomputation ranks
the families the same way (pointer re-walk $>$ replay-rotate-and-swap $>$
re-apply every operator). We report this as compatibility evidence only:
one weak model, no controlled manipulation, no causal claim.

\subsection{Limitations}

Judge coverage is limited to the three implemented families. Each batch is
one model at one temperature with a single run; all batch-2 numbers,
including the decomposition and the recovery gradient, are point estimates
and carry the exploratory mark. The short-chain-to-long-chain extrapolation
question remains unanalysable even in batch 2: per-$k$ sample sizes are
$4$--$30$, neither \texttt{chain-sort} nor \texttt{ret-path} shows a
reliable per-$k$ trend ($R^2 \approx 0$), and the \texttt{stateful-calc}
leave-out test is degenerate (predicting zero against zero). Larger samples
or repeated seeds are required before that question can be answered.

\section{Penalty visibility is a null result, and null is a data point}
\label{sec:q5}

A two-arm experiment (\texttt{blind} vs \texttt{salient}) $\times$ five seeds
$\times$ $100{,}k$ tests whether making the speed penalty visible in the
observation changes the outcome (\texttt{q5/results/train\_config\_q5.json}).
Both arms land at $\mathrm{true} \approx 499.99$ with
\texttt{steps\_alive} $= 500.0$ and $\mathrm{mean\_speed} \approx 0.107$
(\texttt{q5/results/summary\_q5.csv}). The preregistered read-out records the
null and attributes it to option (a)---the penalty has no action surface in
this setting---rather than to a behavioural effect
(\texttt{preregistration\_q5.md:\,L48--L54, L67--L76}). We do not extrapolate
to a general claim that penalty visibility is unimportant; the preregistration
forbids that reading.

\section{Discussion}
\label{sec:disc}

The proxy curve in Section~\ref{sec:decoupling} shows the shape the problem
usually has in practice: the proxy is not merely uninformative, it is
\emph{confidently} informative at the wrong level, and its most confident
level is the one where the true objective is worst. The chained-task audit in
Section~\ref{sec:etl} shows the difference between a measurement that can be
replayed and a measurement that can only be asserted: the same pipeline that
reproduces perfectly also reveals that the model has no partial-credit
behaviour to decompose.

\paragraph{What is next.} One item only, per the reporting discipline: the
chained-task pipeline needs a rerun with a model that is not at floor, since
the two analyses it was built for (error recovery trigger conditions, and
short-to-long chain extrapolation) are blocked by the floor rather than by the
design.

\paragraph{Disclosure.} No capability comparison is made anywhere in this
draft. All agent-side numbers describe one local model at one setting. The
proxy-side numbers describe one environment and one reward family.

\section*{Open items before any submission}
\begin{itemize}
  \item Citations are verified for title/year/venue/DOI but not yet for the
        \texttt{zotero\_key} column, which is unfilled for all 17 rows
        (\texttt{tools/reconcile\_zotero.py} not run).
  \item \texttt{goodhart1975problems} is cited as the 1984 Macmillan reprint
        chapter (\texttt{pp.\ 91--121}); the original 1975 conference paper's
        pagination was not located, and the popular short form
        ``When a measure becomes a target\ldots'' is Strathern 1997, not Goodhart.
  \item The project's own ``three cracks'' taxonomy has no external counterpart:
        the closest work (Manheim \& Garrabrant 2018) uses four categories. The
        three cited laws must not be presented as the source of a three-way split.
  \item The exact wording of the ``decoupling confirmed'' adjudication must
        be copied verbatim from \texttt{docs/HANDOVER.md:\,\S13.11}.
  \item The $\alpha^{*}$ censoring disclosure and the power statement must
        survive into the camera-ready.
\end{itemize}

\nocite{amodei2016concrete,pan2022effects,wei2022chain,ross2011reduction,
kwa2025measuring,schulman2017proximal,raffin2021stable,barto1983neuronlike,
nosek2018preregistration,pineau2021improving,ji2023alignment,morris2024levels}

\bibliographystyle{plain}
\bibliography{refs}

\end{document}
